\chapter{Мастеру}
\label{ch:master}

Глава идёт в порядке реальной работы: от пустого сценария до завершённого подхода. Если вы ведёте игру впервые, прочитайте сначала главу \ref{ch:concepts} --- без неё непонятно, чем сценарий отличается от запущенного мира.

\section{Подготовка сценария}
\label{sec:master:prepare}

\subsection{Создание}
Раздел ``Сценарии'' (\texttt{/scenarios}). Создайте сценарий, дайте название и \textbf{выберите систему правил}.

\begin{quote}
\small
Систему правил потом не поменять: от неё зависят поля всех сущностей. Выбирайте осознанно, см.\ раздел \ref{sec:concepts:rulesystems}.
\end{quote}

\screenshotOptional{user/scenarios.png}{список сценариев}

\subsection{Наполнение}
Внутри сценария (\texttt{/scenarios/<id>}) содержимое разложено по вкладкам, и рядом с названием каждой стоит число --- сколько объектов уже заведено. Это удобный индикатор готовности: пустые вкладки видно сразу.

\screenshotOptional{user/scenario-detail.png}{редактор сценария, вкладка «Сюжет»}

Вкладки: ``Сюжет'' (сюжетные биты), ``Локации'', ``Персонажи'' (заготовки героев игроков), ``NPC'', ``Предметы'', ``Заметки'', ``Счётчики'', ``Todo'' (ваш список задач по подготовке), ``Настройки'' и ``Wiki'' --- свободные текстовые страницы по миру.

У сюжетного бита прямо в карточке видны его экспозиции --- заготовленные куски того, что игроки узнают в этой сцене, с пометками, сколько NPC к ним привязано. Кнопка с глазом открывает предпросмотр, карандаш --- редактирование.

Пошаговый разбор на живом примере --- в главе \ref{ch:tutorials}; справочник по полям --- в главе \ref{ch:objects}.

Полезные привычки:

\begin{itemize}
\item Заполняйте заготовку сцены (что игроки увидят при открытии) сразу, а не во время игры --- за столом на это нет времени.
\item Пишите последствия провала так же подробно, как последствия успеха. Провал --- это поворот сюжета, а не ``ничего не произошло''.
\item Ведите список задач подготовки прямо в сценарии, чтобы не держать в голове.
\end{itemize}

\subsection{Карты и области}
У локации может быть изображение-карта, размеченная областями. Во время игры вы открываете области по одной --- получается туман войны. Разметку делайте заранее: в процессе игры это неудобно.

\subsection{Переиспользование заготовок}
Раздел ``Паки'' (\texttt{/entity-packs}) хранит наборы сущностей, которые можно подключать к разным сценариям --- бестиарий, типовые предметы, готовые NPC. Собрав такой набор один раз, дальше вы экономите время на каждом новом сценарии.

\section{Запуск мира}
\label{sec:master:launch}

Когда сценарий готов, он переводится в запущенное состояние. С этого момента появляется отдельный объект --- запущенный мир, --- и всё, что происходит в игре, меняет его, а не исходную заготовку.

Запущенные миры лежат на \texttt{/launched-scenarios}. Один сценарий можно запустить сколько угодно раз: для разных компаний это будут независимые миры.

\screenshotOptional{user/launched.png}{список запущенных миров}

Страница сама подписывает, что это ``живые копии сценариев между подходами и партиями'', и что исходный сценарий можно править параллельно. Кнопка ``Редактировать'' ведёт в мир, ``Исходник'' --- в заготовку, ``Закрыть'' завершает жизнь мира.

\paragraph{Если планируется серия игр}
Заведите кампанию (\texttt{/campaigns}) до первого подхода. Кампания связывает подходы и переносит между ними состояние персонажей. Без неё каждый подход стартует с начального состояния мира.

\section{Сбор игроков}
\label{sec:master:lobby}

\subsection{Заявки}
Очередь заявок --- на \texttt{/master/applications}. Таблица показывает имя персонажа, систему правил, игрока, статус и когда заявку обновляли; сверху --- сколько заявок ждёт вашего решения, а фильтры по статусу и системе помогают, когда их накопилось много.

\screenshotOptional{user/applications.png}{очередь заявок игроков}

Одобрение создаёт персонажа и привязывает его к игроку.

\subsection{Лобби}
Создайте лобби и разошлите ссылку. В лобби вы видите подключившихся и назначаете каждому персонажа.

\begin{quote}
\small
Назначьте персонажей \textbf{до} запуска игры. Игрок, попавший в сессию без персонажа, увидит пустой экран без объяснений --- он решит, что сайт сломался.
\end{quote}

Лобби можно создать и из Telegram --- бот пришлёт готовую ссылку, см.\ главу \ref{ch:telegram}.

\section{Ведение игры}
\label{sec:master:run}

Экран мастера (\texttt{/session/<id>}) --- это пульт, разделённый на две половины.

\screenshotOptional{user/session-master.png}{экран мастера во время сессии}

\textbf{Слева} --- три вкладки. ``Карта'' показывает текущую локацию с картой и областями, ``Контент'' --- материалы сцены, ``Журнал'' --- хронологию событий. Под ними переключатели ``Локация'', ``Локация сцены'' и ``Таймлайн'', а кнопки ``Локации'' и ``Создать'' открывают выбор места действия.

\textbf{Сверху} --- навигация по сценам: счётчик вида ``1/1'', стрелки вперёд-назад, название текущей сцены и кнопка ``+ Сцена''.

\textbf{Справа} --- блоки, из которых и состоит управление игрой:

\begin{itemize}
\item ``Персонажи'' --- кто в сцене, с портретами и картой локации;
\item ``Время локации'' --- игровые дата и время, которые сдвигаются кнопками $\pm$1д, $\pm$1ч, $\pm$30м, и синхронизируются с сессией кнопками ``К сессии'' и ``В сессию'';
\item ``Сцена (правила)'' --- тип сцены с точки зрения системы правил, например ``Действие'';
\item ``Публичные элементы'' --- то, что игроки видят;
\item ``Приватные элементы'' --- то, что видите только вы.
\end{itemize}

Деление на публичные и приватные элементы --- главный рычаг мастера: перенос объекта из приватных в публичные и есть тот момент, когда игроки что-то узнают.

Кнопка ``Настройки'' в правом нижнем углу открывает параметры сессии.

\subsection{Сцены}
Открытая сцена --- то, что игроки сейчас видят. Переключение сцены мгновенно меняет картинку у всех. Заготовленный текст сцены уходит игрокам в момент открытия.

\subsection{Сюжетные биты и таймлайн}
Сюжетные биты, заготовленные в редакторе сценария, служат вашим ориентиром по ходу истории; игрокам сама их структура не видна. Хронология уже случившегося доступна во вкладке ``Таймлайн'' слева и в ``Журнале''.

\subsection{Карта}
Открывайте области карты по мере продвижения группы. Закрытые области игрокам не видны.

\subsection{Показать объект}
Любой объект --- NPC, предмет, локацию --- можно вывести крупно у всех участников. Самый быстрый способ представить нового персонажа или показать находку.

\subsection{Счётчики}
Часы, ресурсы, шкалы угрозы. Заводятся на вкладке ``Счётчики'' в редакторе сценария, а во время игры попадают в элементы сцены. Пока счётчик лежит в приватных элементах, его видите только вы; перенесённый в публичные --- он виден игрокам, и его рост хорошо работает на нарастающее напряжение.

\subsection{Заметки и рассылки}
Три разных вещи:

\begin{itemize}
\item личные заметки --- видите только вы;
\item заметки для игроков --- видны всем;
\item адресная рассылка --- уходит конкретному игроку, остальные её не видят. Игрок может ответить.
\end{itemize}

\subsection{Музыка}
Раздел ``Аудио'' (\texttt{/audio}) и очередь треков в сессии. Воспроизведение синхронизировано: трек играет у всех участников одновременно.

\subsection{Броски и действия}
Действия игроков проходят по шагам (см.\ раздел \ref{sec:player:actions}). Вы участвуете там, где нужен выбор мастера: определить последствие, назначить урон, подтвердить изменения.

Отменить чужое действие вы можете на любом шаге, в том числе после броска. У игрока такой возможности после броска нет.

\subsection{Зрители}
Если игру кто-то смотрит --- создайте комнату наблюдателя и раздайте ссылку, см.\ главу \ref{ch:observer}.

\section{Завершение подхода}
\label{sec:master:finish}

Подход завершается вручную.

\begin{quote}
\small
\textbf{Действие необратимо.} Завершённый подход не открыть заново, а интерфейс не показывает заранее, что именно перенесётся в следующий. Убедитесь, что все изменения внесены, прежде чем нажимать.
\end{quote}

Если подход входит в кампанию, состояние персонажей, инвентарь и счётчики переедут в следующий подход. Мир остаётся жив --- следующую встречу вы начинаете из него, а не с нуля.

\section{Что стоит знать заранее}
\label{sec:master:caveats}

Несколько особенностей текущей версии, о которые спотыкаются чаще всего:

\begin{itemize}
\item \textbf{Правки локаций во время сессии могут не сохраниться.} Менять состав локаций надёжнее в редакторе сценария, а не на ходу.
\item \textbf{Полностью проработана только Dungeon World.} На других системах ведите броски за столом вручную.
\item \textbf{Ссылку на комнату наблюдателя нельзя отозвать.} Кому дали --- тот сможет смотреть и дальше.
\end{itemize}
